\documentclass[tikz,border=5mm]{standalone}
\usepackage{xcolor}

% Farbdefinition für den roten Kreisbogen
\definecolor{crimsonred}{RGB}{186, 45, 63}

\begin{document}
\begin{tikzpicture}[scale=1.5]
    % Parameter für Radius und Winkel
    \def\r{2.2}
    \def\angle{115}
    
    % Zeichne den vollständigen Kreis im Hintergrund (hellgrau und dünn)
    \draw[gray!30, thin] (0,0) circle (\r);
    
    % Zeichne die beiden Radien
    \draw[gray!60, semithick] (0,0) -- (\r,0);
    \draw[gray!60, semithick] (0,0) -- (\angle:\r);
    
    % Winkelbogen und Beschriftung für phi
    \draw[black, thin] (0:0.45) arc (0:\angle:0.45);
    \node at ({\angle/2}:0.25) {$\varphi$};
    
    % Hervorgehobener Kreisbogen b
    \draw[crimsonred, thick] (0:\r) arc (0:\angle:\r);
    
    % Beschriftung der Radien und des Bogens
    \node[above] at (0.5*\r, 0.02) {$r$};
    \node[right, xshift=2pt] at (\angle:0.5*\r) {$r$};
    \node[above right] at (60:\r) {$b$};
    
    % Kleine weiße Kreise an den Knotenpunkten zeichnen
    \filldraw[fill=white, draw=black, thin] (0,0) circle (1.5pt);
    \filldraw[fill=white, draw=black, thin] (\r,0) circle (1.5pt);
    \filldraw[fill=white, draw=black, thin] (\angle:\r) circle (1.5pt);
    
\end{tikzpicture}
\end{document}
